% Fichier complété par tools/complete_missing_corrections.py.
% Source : ELEC/ELEC-01/536_CircuitElec/536_CircuitElec.tex
% Statut : rédigé (1 question(s)).
\begin{corrigebox}[Corrigé rédigé]
\CorrigeQuestion{1}
On repère les nœuds puis on choisit le potentiel du rail inférieur
            nul et celui du générateur égal à $E$. Pour chaque nœud $k$, la loi des nœuds
            donne
            \[\sum_{j}\frac{V_k-V_j}{R_{kj}}=0.\]
            La résolution du système linéaire fournit tous les potentiels. Pour chaque
            résistance $R_i$ entre les nœuds $a$ et $b$ :
            $U_i=V_a-V_b$ et $I_i=U_i/R_i$. Les courants orientés dans le sens opposé à la
            valeur calculée sont négatifs. Cette méthode s'applique sans ambiguïté au
            réseau de la figure et constitue l'expression demandée en fonction de $E$ et
            des $R_i$.
\end{corrigebox}
